\chapter{Operations and runbooks}\label{ch:ops}
\chapterstatus{draft}{2026-11-21}{\src{infra/aws}, \src{infra/aws/deploy.sh}}

\section{Access}
Operators reach \code{slotty-prod} from the management account; nobody has a user in it.
\begin{description}
  \item[Console] Sign in to the management account, then \emph{Switch role} to account
    \code{slotty-prod} with role \code{OrganizationAccountAccessRole}.
  \item[CLI] A profile that assumes the same role; every command below uses it.
\end{description}
\begin{lstlisting}
[profile slotty-prod]
role_arn       = arn:aws:iam::<slotty_account_id>:role/OrganizationAccountAccessRole
source_profile = default
region         = us-east-1
\end{lstlisting}
The values the commands need (subnets, security group, cluster) are in
\code{terraform -chdir=infra/aws/app output github_actions_variables}.

\section{Runbooks}
\subsection{Deploy}
Publish a release (\cref{sec:releases}); merging alone deploys nothing. To redeploy without a code
change (for example after changing configuration in Terraform), run the \emph{Deploy} workflow by
hand from the running release's tag; it reuses the images and skips the migrate task unless
\emph{force\_migrate} is ticked. \code{curl https://<domain>/health} shows the release that is
serving (\code{version}, the tag).

\subsection{Hotfix}
Branch off \code{master}, fix, merge, then publish the next patch release (\code{v1.2.1}). A
release is always a commit on \code{master}, so if \code{master} already holds unreleased work
that mustn't ship yet, revert it first and merge it again after the release.

\subsection{A deploy failed}
The rollout job's log says which step stopped it and, for a service, prints its recent ECS events.
\begin{description}
  \item[Migrations failed] The services were not touched. Fix the migration in a new pull request.
  \item[A service did not roll out] The circuit breaker rolled it back, or its tasks didn't become
    healthy within 10 minutes. The service events and the task's log group (\code{/slotty/web} or
    \code{/slotty/worker}) say why; usually a missing setting or a crash at start-up.
  \item[\route{/health} doesn't report the release] The tasks run but the public path doesn't
    reach them, or the database check fails: see \route{/health}'s answer and the API Gateway
    log group.
  \item[A newer release is already live] Not a failure: an older release's run was queued behind
    a newer one. To go back on purpose, see below.
  \item[package.json says … but the tag is …] The release was tagged before the
    \emph{Release vX.Y.Z} pull request merged, or on the wrong commit. Delete the release and its
    tag (admins only) and publish again.
\end{description}

\subsection{Roll back}
Run \emph{Deploy} by hand from the older release's tag with \emph{rollback} ticked. Its images are
still in ECR, so nothing is rebuilt; the migrate task doesn't run and \route{/health} confirms the
older release. If GitHub is unavailable, point the services at the previous task-definition
revisions instead:
\begin{lstlisting}[language=bash]
aws ecs update-service --profile slotty-prod --cluster slotty --service web    --task-definition slotty-web:<n>
aws ecs update-service --profile slotty-prod --cluster slotty --service worker --task-definition slotty-worker:<n>
\end{lstlisting}
Migrations are not rolled back automatically. Write them so the previous release still works
against the new schema (expand first, contract in a later release).

\subsection{Run a one-off command}
Run it as a task with the worker's definition and its secrets. Creating an admin:
\begin{lstlisting}[language=bash]
aws ecs run-task --profile slotty-prod --cluster slotty --task-definition slotty-worker \
  --launch-type FARGATE --network-configuration "awsvpcConfiguration={subnets=[<a>,<b>],securityGroups=[<sg>],assignPublicIp=DISABLED}" \
  --overrides '{"containerOverrides":[{"name":"worker","command":["node","dist/jobs/create-admin.js"],
    "environment":[{"name":"ADMIN_EMAIL","value":"…"},{"name":"ADMIN_NAME","value":"…"},{"name":"ADMIN_PASSWORD","value":"…"}]}]}'
\end{lstlisting}
\begin{risk}
  Overrides are visible in the task's details and in CloudTrail. Treat a password passed this way
  as temporary and change it after the first sign-in.
\end{risk}

\subsection{Give someone platform access}\label{sec:ops-platform}
Usage insights and feedback (\route{/platform}) need a dedicated account: not an admin and not in
any course. Grant it as a one-off command, like creating an admin, with
\code{["node","dist/jobs/platform-access.js","grant"]} and \envvar{PLATFORM_EMAIL},
\envvar{PLATFORM_NAME}, \envvar{PLATFORM_PASSWORD} (locally: \code{bun run platform:grant}). An
existing account keeps its password; a new one is created active. \code{revoke} (needs only
\envvar{PLATFORM_EMAIL}) removes access and signs the account out. Both are audited as
\code{platform.grant} / \code{platform.revoke}, which only the database shows. There is
deliberately no button for this anywhere in the app.

\subsection{Reach the database}
Port-forward through the NAT instance with Session Manager (no SSH, no bastion), then connect
with the credentials in \code{/slotty/DATABASE_URL}:
\begin{lstlisting}[language=bash]
aws ssm start-session --profile slotty-prod --target <nat_instance_id> \
  --document-name AWS-StartPortForwardingSessionToRemoteHost \
  --parameters host=<db_address>,portNumber=5432,localPortNumber=15432
\end{lstlisting}

\subsection{Someone can't sign in: ``Too many failed sign-ins''}
The account (or their network) hit the sign-in limit (\cref{ch:rules}, ADR-026). The wait ends by
itself within 15 minutes. To lift it at once, send them a reset from \route{/admin/users/[id]}
(\ui{Send password reset}): resetting through the link clears the account's failures. A whole
campus locked out at once means one address made 50 failed sign-ins in 15 minutes; it also
clears by itself. Nobody needs database access for either.

\subsection{Email is not arriving}
Check \route{/admin/email}: failed and waiting messages show the SMTP error. A permission or
credential error affects every message (\code{slotty-worker-failures} fires); a bounce or
complaint arrives as an SNS email and the address is suppressed. Retry from the console after
fixing the cause. Rotate the SMTP credentials with
\code{terraform apply -replace=aws_iam_access_key.ses_smtp}, then redeploy.

\subsection{Support mail is not arriving}
\code{support@slotty.msaeedsaeedi.com} is forwarded by Cloudflare Email Routing (ADR-021). Two
places to check. DNS: the domain's MX records must point to \code{route1/2/3.mx.cloudflare.net}
(\src{infra/aws/app/dns.tf}; look them up with any DNS tool). Cloudflare: \ui{Email}
\textrightarrow{} \ui{Email Routing} on \code{msaeedsaeedi.com} must have the rule for the
address and a verified destination; its activity log shows each message as forwarded, dropped or
rejected. Cloudflare's dashboard may report the subdomain's records as missing, because they live
in Route 53 rather than in its own zone; that is expected.

\subsection{Usage data looks wrong or empty}
\route{/platform/insights} says so when \envvar{ANALYTICS_SECRET} isn't set on the web tasks (check
the task definition's secrets). Raw events appear within about 10 seconds of a visit; the
13-month trend fills from the next hourly worker run (look for \code{analyticsRolled=} in
\code{/slotty/worker}; an \code{analytics failed} line names the error). The ingest endpoint logs
\code{[analytics] could not record events} in \code{/slotty/web} and still answers 204, so the
app keeps working. Rotating the secret
(\code{terraform apply -replace=random_password.analytics}, then a deploy) makes every person
count as new from then on. No alarm is spent on analytics: it is not part of the service.

\subsection{An alarm fired}
\begin{xltabular}{\linewidth}{@{}P{40mm} L@{}}
  \toprule
  \thead{Alarm} & \thead{First steps}\\
  \midrule
  \code{web-down}, \code{worker-down} & ECS console: the service's events and stopped tasks' reasons. Spot capacity, a failing health check and a bad image are the usual causes.\\
  \code{5xx} & \code{/slotty/api} access log for the failing paths, then \code{/slotty/web} for the errors.\\
  \code{worker-failures} & \code{/slotty/worker}: a job's error, or repeated send failures (see above).\\
  \code{db-*} & RDS Performance and the slow-statement log; consider \tfvar{db_instance_class}.\\
  \code{nat-status-failed} & EC2 automatic recovery usually restarts it; otherwise replace it with \code{terraform apply -replace=aws_instance.nat[0]}.\\
  \bottomrule
\end{xltabular}

\begin{todo}
  \begin{itemize}
    \item Restore drill: restore a point-in-time copy of the database to a temporary instance,
      check it, delete it. Record the date and the time it took.
    \item Incident process: who is alerted, severity levels, and the post-incident review.
    \item Scaling playbook for a large demo day: when to raise \tfvar{web_min} beforehand.
  \end{itemize}
\end{todo}
